\section*{Hoofdstuk \ref{ch:b3:surfaces-catalysis} --- Oppervlakken, adsorptie en heterogene katalyse}

\begin{solution}{exo:b3:surfaces-catalysis:1}
$K = \theta/((1 - \theta)p) = 0.40/(0.60 \times 2.0) = \qty{0.33}{kPa^{-1}}$; bij \qty{10}{kPa} is $\theta = 3.33/4.33 = 0.77$.
\end{solution}

\begin{solution}{exo:b3:surfaces-catalysis:2}
$-20$: \omterm{def:b3:surfaces-catalysis:physi-chemi}{fysisorptie} (de grootteorde van een condensatie-enthalpie); $-150$:
\omterm{def:b3:surfaces-catalysis:physi-chemi}{chemisorptie}, die als enige dissociatief kan zijn.
\end{solution}

\begin{solution}{exo:b3:surfaces-catalysis:3}
(100): één topplaats, twee brugplaatsen (vier bindingen, elk gedeeld door
twee atomen), één viervoudige holte (vier holten, elk gedeeld door vier
atomen). (111): één topplaats, drie brugplaatsen, twee drievoudige holten
(zes driehoeken rond elk atoom, elk gedeeld door drie).
\end{solution}

\begin{solution}{exo:b3:surfaces-catalysis:4}
$\num{3.0e-6}/\num{1.5e-5} = \qty{0.20}{s^{-1}}$.
\end{solution}

\begin{solution}{exo:b3:surfaces-catalysis:5}
$\theta/(1 - \theta) = (Kp)^{1/2}$: $0.111$ bij \qty{1.0}{kPa}, verdubbeld bij \qty{4.0}{kPa}: $0.222$, dus $\theta = 0.18$.
\end{solution}

\begin{solution}{exo:b3:surfaces-catalysis:6}
Noemer $1 + 10 + 2.5 = 13.5$: $\theta_{\ce{CO}} = 0.74$, $\theta_{\ce{O2}} = 0.19$, vrije plaatsen 0.07. De
snelheid volgens Langmuir-Hinshelwood, evenredig met $\theta_{\ce{CO}}\theta_{\ce{O2}}$, wordt beperkt door de
schaarste aan zuurstof op een oppervlak dat vol CO zit.
\end{solution}

\begin{solution}{exo:b3:surfaces-catalysis:7}
$\Delta_{\mathrm{ads}}H = -R\ln(8.0/1.0)/(1/300 - 1/350) = -8.314 \times 2.079/\num{4.76e-4} = \qty{-36}{kJ/mol}$.
\end{solution}

\begin{solution}{exo:b3:surfaces-catalysis:8}
$n_{\mathrm m} = 120/22\,414 = \qty{5.35e-3}{mol}$; oppervlakte $= \num{5.35e-3} \times \num{6.022e23} \times \num{0.162e-18} =
\qty{522}{m^2.g^{-1}}$.
\end{solution}

\begin{solution}{exo:b3:surfaces-catalysis:9}
Langmuir-Hinshelwood: beide beginstoffen moeten op naburige plaatsen
geadsorbeerd zijn; CO bindt sterk en laat bij hoge druk geen ruimte over
voor zuurstof.
\end{solution}

\begin{solution}{exo:b3:surfaces-catalysis:10}
Lage \omterm{def:b3:surfaces-catalysis:coverage}{bedekkingsgraad}: $100 - 60 = \qty{40}{kJ/mol}$; hoge \omterm{def:b3:surfaces-catalysis:coverage}{bedekkingsgraad}: \qty{100}{kJ/mol}. De grafiek van $\ln r$
tegen $1/T$ is steil (helling $-100/R$) bij lage temperatuur, waar het
oppervlak vol is, en vlakker (helling $-40/R$) bij hoge temperatuur, waar het leegloopt: een kromme
die tussen beide ombuigt.
\end{solution}

\begin{solution}{exo:b3:surfaces-catalysis:11}
$4 \times 0.10 = 0.40$ van de plaatsen is geblokkeerd: er blijft 60\,\% van de
activiteit over voor één plaats, ongeveer $0.60^2 = 36\,\%$ voor twee naburige vrije
plaatsen.
\end{solution}

\begin{solution}{exo:b3:surfaces-catalysis:12}
Het middelste (\qty{-250}{kJ/mol}): het eerste bindt stikstof te zwak om \ce{N2}
gemakkelijk te verbreken, het derde te sterk, zodat de stikstofatomen op
het oppervlak blijven en het blokkeren.
\end{solution}

\begin{solution}{pb:b3:surfaces-catalysis:1}
\textbf{1.} Stikstof condenseert bij atmosferische druk rond \qty{77}{K}: het
volledige gebied van relatieve drukken is bereikbaar en er vormen zich
meerlagen.
\textbf{2.} Onder 0.05 schendt de heterogeniteit van het oppervlak, boven
0.30 de condensatie in de poriën, de aannamen van BET.
\textbf{3.} \num{1.278e-3}, \num{2.410e-3}, \num{3.453e-3}, \num{4.621e-3}, \num{5.659e-3}, \num{6.813e-3} \unit{g.cm^{-3}}.
\textbf{4.} $v_{\mathrm m} = 1/(0.02205 + 0.000180) = \qty{45.0}{cm^3.g^{-1}}$.
\textbf{5.} $c = 1 + 0.02205/0.000180 = 124$: de eerste laag bindt veel sterker
dan de vloeistof.
\textbf{6.} Ze is zeer klein vergeleken met de helling maal $x$, dus haar
relatieve fout is groot; maar $v_{\mathrm m}$ hangt af van de som van helling en afsnijding, en die wordt door de helling
gedomineerd.
\textbf{7.} $45.0/22\,414 = \qty{2.01e-3}{mol.g^{-1}}$.
\textbf{8.} $\num{2.01e-3} \times \num{6.022e23} \times \num{0.162e-18} = \qty{196}{m^2.g^{-1}}$.
\textbf{9.} De aluminiumoxidedrager: het blootliggende platina (deel III)
heeft ongeveer \qty{0.9}{m^2.g^{-1}}.
\textbf{10.} Poriën van moleculaire breedte vullen zich volledig bij zeer
lage druk, niet laag na laag, zodat er geen \omterm{def:b3:surfaces-catalysis:coverage}{monolaag} te herkennen is.
\textbf{11.} $2 \times 0.205/22\,414 = \qty{1.83e-5}{mol.g^{-1}}$.
\textbf{12.} $0.0100/195.08 = \qty{5.13e-5}{mol.g^{-1}}$.
\textbf{13.} $D = 1.83/5.13 = 0.357$.
\textbf{14.} $a^3/4 = (0.3923)^3/4 = \qty{0.01509}{nm^3}$.
\textbf{15.} $\frac{\sqrt3}{4}(0.3923)^2 = \qty{0.0666}{nm^2}$, het dichtste vlak; (100) en (110) geven 0.0770 en
\qty{0.1088}{nm^2}, en het gemiddelde van de drie plaatsdichtheden geeft \qty{0.0807}{nm^2}.
\textbf{16.} Een bol bevat $\pi d^3/6v_{\mathrm{at}}$ atomen, waarvan er $\pi d^2/a_{\mathrm{at}}$ aan het oppervlak liggen:
$D = 6v_{\mathrm{at}}/(a_{\mathrm{at}}d)$.
\textbf{17.} $d = 6 \times 0.01509/(0.0807 \times 0.357) = \qty{3.1}{nm}$.
\textbf{18.} Eén H per Pt-atoom aan het oppervlak; bolvormige deeltjes van
één grootte (het resultaat is een naar oppervlakte gewogen gemiddelde); al
het platina gereduceerd en toegankelijk; geen waterstof die naar de
drager overloopt.
\textbf{19.} $\num{2.0e-5}/\num{1.83e-5} = \qty{1.1}{s^{-1}}$.
\textbf{20.} $\num{2.0e-5}/0.0100 = \qty{2.0e-3}{mol.g^{-1}.s^{-1}}$ per gram platina.
\textbf{21.} $D = 1.12/10 = 0.112$ in plaats van 0.357: een factor 3.2.
\textbf{22.} Dat de reactie structuurgevoelig is: haar actieve plaatsen
(randen, hoeken, bepaalde vlakken) vormen geen constante fractie van het
oppervlak.
\textbf{23.} Ze telt de gebeurtenissen per \omterm{def:b3:complex-kinetics:enzyme}{actieve plaats} en neemt zo het
effect weg van de hoeveelheid metaal die blootligt.
\textbf{24.} \textbf{Gemiddelde diameter van de platinadeeltjes $\approx \qty{3.1}{nm}$.}
\end{solution}
